% ============================================================
%  Глава 4 — Pure Pursuit: алгоритм следования по пути
%  Файл: pi_nodes/nodes/path_recorder_node.py
% ============================================================
\chapter{Pure Pursuit: алгоритм возврата домой}
\label{ch:pure_pursuit}

\begin{filebox}[title={Файл исходного кода}]
\filepath{pi\_nodes/nodes/path\_recorder\_node.py} \quad (342 строки, Python)
\end{filebox}

\section{Постановка задачи}

Модуль записывает путь робота в виде набора опорных точек и затем проводит
робота обратно по тому же пути в обратном порядке (возврат домой).

\begin{itemize}
  \item \textbf{Запись}: каждые $3\,\text{см}$ или $4.5^\circ$ поворота.
  \item \textbf{Воспроизведение}: алгоритм Pure Pursuit с финальным фазой
    точного наведения.
\end{itemize}

\section{Запись путевых точек}

\begin{filebox}[title={\filepath{path\_recorder\_node.py}, строки 135--148 --- \_maybe\_record\_waypoint()}]
\begin{lstlisting}[style=pythonstyle, numbers=left, firstnumber=139]
dx = self._x - self._last_record_x
dy = self._y - self._last_record_y
dist = math.sqrt(dx * dx + dy * dy)
dtheta = abs(self._angle_diff(self._theta, self._last_record_theta))

if dist >= RECORD_INTERVAL_M or dtheta >= RECORD_INTERVAL_RAD:
    self._path.append((self._x, self._y, self._theta))
    # RECORD_INTERVAL_M = 0.03 m, RECORD_INTERVAL_RAD = 0.08 rad
\end{lstlisting}
\end{filebox}

\textbf{Расстояние между точками} (Евклидова метрика):

\begin{equation}
  d = \sqrt{(x - x_{\text{last}})^2 + (y - y_{\text{last}})^2}
  \label{eq:waypoint_dist}
\end{equation}

Запись происходит при: $d \geq 0.03\,\text{м}$ ИЛИ $|\Delta\theta| \geq 0.08\,\text{рад}$.

\section{Алгоритм Pure Pursuit}
\label{sec:pure_pursuit_algo}

\begin{filebox}[title={\filepath{path\_recorder\_node.py}, строки 206--243 --- \_control\_loop()}]
\begin{lstlisting}[style=pythonstyle, numbers=left, firstnumber=206]
# Pure pursuit: simultaneous linear + angular
bearing = math.atan2(dy, dx)
angle_error = self._angle_diff(bearing, self._theta)

if in_final_approach:
    angular   = FINAL_ANGULAR_KP * angle_error    # Kp = 3.0
    angular   = max(-REPLAY_ANGULAR_SPEED, min(REPLAY_ANGULAR_SPEED, angular))
    cos_factor = max(0.0, math.cos(angle_error))
    dist_ratio = dist / FINAL_APPROACH_RADIUS      # 0..1
    linear     = FINAL_LINEAR_SPEED * cos_factor * dist_ratio
    linear     = max(REPLAY_MIN_LINEAR * 0.6 * cos_factor, linear)
    if abs(angle_error) > 0.8:
        linear = 0.0
else:
    angular    = REPLAY_ANGULAR_KP * angle_error  # Kp = 2.5
    angular    = max(-REPLAY_ANGULAR_SPEED, min(REPLAY_ANGULAR_SPEED, angular))
    cos_factor = max(0.0, math.cos(angle_error))
    dist_factor = min(1.0, dist / 0.15)
    linear     = REPLAY_LINEAR_SPEED * cos_factor * dist_factor
    linear     = max(REPLAY_MIN_LINEAR * cos_factor, linear)
    if abs(angle_error) > 1.2:
        linear = 0.0
\end{lstlisting}
\end{filebox}

\subsection{Вычисление угла на цель}

Пеленг (bearing) на текущую целевую точку $(t_x, t_y)$:

\begin{equation}
  \beta = \mathrm{atan2}(t_y - y,\; t_x - x)
  \label{eq:bearing}
\end{equation}

\textbf{Угловая ошибка} (кратчайший путь в $[-\pi, \pi]$):

\begin{equation}
  e_\theta = \mathrm{angle\_diff}(\beta, \psi) \in [-\pi, \pi]
  \label{eq:angle_error}
\end{equation}

\subsection{Пропорциональный угловой регулятор}

\begin{equation}
  \omega = K_P \cdot e_\theta, \quad
  \omega \in [-\omega_{\max}, +\omega_{\max}]
  \label{eq:angular_control}
\end{equation}

\begin{center}
\begin{tabular}{lll}
\toprule
\textbf{Режим} & $K_P$ & $\omega_{\max}$\\
\midrule
Обычный Pure Pursuit & $2.5$ & $0.8\,\text{рад/с}$\\
Финальный подход & $3.0$ & $0.8\,\text{рад/с}$\\
\bottomrule
\end{tabular}
\end{center}

\subsection{Косинусное масштабирование линейной скорости}

\textbf{Ключевая идея Pure Pursuit}: при большой угловой ошибке снижать
линейную скорость через косинус:

\begin{equation}
  \cos\text{-factor} = \max\bigl(0,\, \cos(e_\theta)\bigr)
  \label{eq:cos_factor}
\end{equation}

\begin{equation}
  v = v_{\max} \cdot \cos\text{-factor} \cdot d_{\text{factor}}
  \label{eq:linear_speed}
\end{equation}

Коэффициент дистанции:

\begin{equation}
  d_{\text{factor}} = \min\!\left(1,\, \frac{d}{0.15\,\text{м}}\right)
  \label{eq:dist_factor}
\end{equation}

При $|e_\theta| > \pi/2$: $\cos(e_\theta) \leq 0 \Rightarrow v = 0$
(только поворот на месте).

\subsection{Финальный подход}

В радиусе $0.12\,\text{м}$ от конечной точки:

\begin{equation}
  v = v_{\text{final}} \cdot \cos(e_\theta) \cdot \frac{d}{r_{\text{final}}}
  \label{eq:final_approach}
\end{equation}

Скорость пропорциональна оставшемуся расстоянию --- плавное торможение.

\begin{center}
\begin{tabular}{lll}
\toprule
\textbf{Параметр} & \textbf{Значение} & \textbf{Описание}\\
\midrule
\texttt{REPLAY\_LINEAR\_SPEED}   & $0.20\,\text{м/с}$ & Скорость воспроизведения\\
\texttt{FINAL\_LINEAR\_SPEED}    & $0.08\,\text{м/с}$ & Скорость финального подхода\\
\texttt{REPLAY\_LOOKAHEAD\_M}    & $0.10\,\text{м}$   & Дистанция look-ahead\\
\texttt{REPLAY\_HOME\_TOLERANCE} & $0.015\,\text{м}$  & Допуск финальной точки\\
\texttt{REPLAY\_GOAL\_TOLERANCE} & $0.04\,\text{м}$   & Допуск промежуточных точек\\
\bottomrule
\end{tabular}
\end{center}

\section{Нормализация угла}

\begin{filebox}[title={\filepath{path\_recorder\_node.py}, строки 312--320 --- \_angle\_diff()}]
\begin{lstlisting}[style=pythonstyle, numbers=left, firstnumber=313]
@staticmethod
def _angle_diff(target, current):
    """Shortest signed angle difference, in [-pi, pi]."""
    d = target - current
    while d > math.pi:
        d -= 2.0 * math.pi
    while d < -math.pi:
        d += 2.0 * math.pi
    return d
\end{lstlisting}
\end{filebox}

\begin{equation}
  \Delta\theta = \bigl[(\beta - \psi + \pi) \bmod 2\pi\bigr] - \pi
  \label{eq:angle_normalize}
\end{equation}

Это гарантирует, что угловая ошибка всегда в диапазоне $[-\pi, \pi]$,
а поворот всегда происходит по кратчайшему пути.

\section{Геометрия алгоритма Pure Pursuit}

\begin{center}
\begin{tikzpicture}[scale=1.4]
  % Robot
  \filldraw[blue!70] (0,0) circle (0.12);
  \draw[-{Stealth}, blue!70, thick] (0,0) -- (1.0, 0)
    node[right, font=\small]{$\psi$ (курс)};
  \node[below, font=\small, blue!70] at (0,-0.15) {Робот};

  % Target point
  \filldraw[red!70] (1.5, 1.2) circle (0.08);
  \node[right, font=\small, red!70] at (1.58, 1.2) {Цель $(t_x, t_y)$};

  % Bearing line
  \draw[-{Stealth}, red!50, dashed, thick] (0,0) -- (1.5, 1.2)
    node[midway, above left, font=\small]{$d$};

  % Angle annotation
  \draw[->] (0.5, 0) arc[start angle=0, end angle=38.7, radius=0.5];
  \node[font=\small] at (0.72, 0.2) {$e_\theta$};

  % Bearing angle label
  \draw[->] (0.3, 0) arc[start angle=0, end angle=38.7, radius=0.3];
  \node[font=\tiny, gray] at (0.48, 0.08) {$\beta$};
\end{tikzpicture}
\end{center}

$$\beta = \mathrm{atan2}(t_y - y, t_x - x), \quad e_\theta = \beta - \psi$$

\section{Теоретическое обоснование: замкнутая система управления}

\subsection{Структура контура управления}

Алгоритм Pure Pursuit реализует \textbf{замкнутый P-регулятор}
(гл.~\ref{ch:theory}, раздел~\ref{sec:p_controller}) для угловой
скорости:

\begin{equation}
  \omega = K_P \cdot e_\theta, \quad e_\theta = \beta - \psi
  \label{eq:pp_closed_loop}
\end{equation}

Передаточная функция замкнутой системы <<регулятор + робот>>
(робот как интегратор $\dot{\psi} = \omega$):

\begin{equation}
  W_{\text{зс}}(p) = \frac{K_P / p}{1 + K_P / p} = \frac{K_P}{p + K_P}
  \label{eq:pp_tf}
\end{equation}

Полюс замкнутой системы: $\lambda = -K_P = -2.5$,
что даёт экспоненциальную сходимость $e_\theta(t) = e_\theta(0)\,e^{-2.5t}$
с постоянной времени $\tau = 1/K_P = 0.4\,\text{с}$.

\subsection{Нелинейное масштабирование и устойчивость}

Косинусное масштабирование $v = v_{\max}\cos(e_\theta)$ создаёт
\textbf{нелинейную связь} между угловой ошибкой и линейной скоростью.
Это эквивалентно функции Ляпунова:

\begin{equation}
  V(e_\theta) = 1 - \cos(e_\theta) \geq 0
  \label{eq:pp_lyapunov}
\end{equation}

При $|e_\theta| < \pi/2$: $\dot{V} = \sin(e_\theta)\cdot\dot{e}_\theta < 0$
(при правильном знаке $K_P$), что гарантирует
\textbf{асимптотическую устойчивость} по Ляпунову.

\subsection{Связь с кинематической моделью}

Полная нелинейная модель робота с Pure Pursuit
(модель unicycle из гл.~\ref{ch:odometry}):

\begin{equation}
  \begin{cases}
    \dot{x} = v\cos\theta \cdot \cos(e_\theta) \cdot d_{\text{factor}} \\
    \dot{y} = v\sin\theta \cdot \cos(e_\theta) \cdot d_{\text{factor}} \\
    \dot{\theta} = K_P \cdot e_\theta
  \end{cases}
  \label{eq:pp_full_model}
\end{equation}

Это нелинейная система вида $\dot{\vect{x}} = \vect{f}(\vect{x})$
(гл.~\ref{ch:theory}, формула~\eqref{eq:nonlinear_system}),
линеаризация которой в окрестности $e_\theta = 0$ даёт
устойчивую линейную систему с полюсом $\lambda = -K_P$.
